Este estudo aborda a otimização operacional de sistemas de abastecimento de água, com o objetivo de minimizar os custos e melhorar a eficiência no uso de energia elétrica. A fundamentação teórica é baseada em análises temporais do consumo de água, padrões de utilização, correlações estatísticas e técnicas de otimização. O modelo com variáveis binárias adotado leva em conta tarifas de energia que mudam ao longo do tempo e políticas de consumo, visando desenvolver estratégias operacionais para diminuir os custos com energia. A análise dos dados evidenciou padrões que podem ser utilizados para a formulação de estratégias eficazes de gestão. A otimização implementada indicou a possibilidade de economias significativas. Planos futuros incluem a ampliação da base de dados, análises de sensibilidade e o desenvolvimento de novas políticas de gestão, a fim de promover um gerenciamento mais eficiente da energia elétrica.

% Separe as palavras-chave por ponto
\palavraschave{Gestão de Recursos Hídricos; Eficiência Energética; Otimização de Custos; Modelagem Matemática.}